% !TEX root = ../monograph.tex
\subsection{Intervals of scales crossing changes of stage}
\label{ch17:variable-windows}

The finite comparison in Proposition~\ref{ch11:n-mark-comparison}
extends to the intervals of radii just described. We give the estimates
because requiring every interval of scales to lie inside a plateau would leave
gaps in radius coverage.

Fix the maximum scaling dimension of the insertions, their number, and
the orders of the derivatives to be estimated. In this paragraph
$C_j$ includes the linear column constants with insertions, the curvature-response
row constants and the finite polynomial-source propagation constants
for this request, all fixed \emph{before} the block ratio of stage $j$.
Write $\mathscr L_j$ for the ratio chosen at stage $j$, so
$L_i=\mathscr L_{J_i}$ with $J_i=J(i+1)$ in the scale convention
of Chapter~\ref{ch10:trajectory}. The ratio choices may require
\[
 \varepsilon_j:=\frac{\log(2+C_j)}{\log\mathscr L_j}\longrightarrow0.
\]
For example one may impose $\varepsilon_j\le j^{-2}$.
Constants depending on the already selected ratio, such as nonlinear
jet bounds of the critical configuration, are not included here. They enter the local
small errors and endpoint constants, and are handled by sufficiently large lower bounds on the reference coupling.
These local errors are summed after transfer; their ratio-dependent
constants are never multiplied once per step. Put $H_n=\exp((\log\log\widehat\beta_n)^2)$. Only a bounded number
of endpoint factors is multiplied, giving $H_n^{o(1)}$.
This distinction avoids a circular ratio choice.

For $b$ in a fixed compact interval of sufficiently large positive
numbers, move backwards from scale $n$, adding $\log L_i$ at each step.
Stop when this sum first reaches $(b+2)\log H_n$. Starting from that
scale, continue until the additional sum first reaches $10\log B_n$. Here $B_n$ is the radius of the restarted representation;
$\log B_n=O(\log\log\widehat\beta_n)$.
The schedule bound $L_i\le\log\widehat\beta_i$ makes each
excess over the target $O(\log\log\widehat\beta_n)=o(\log H_n)$.
The entire interval of scales has logarithmic length $O(\log H_n)$, and the
number of steps is $O(\log H_n)$ since every $L_i\ge3$.
It is therefore shorter than $(\log\log\widehat\beta_n)^4$,
as required by Proposition~\ref{ch11:restart-family}.
All its scale indices tend to infinity. The increment estimate is a
constant times the logarithmic block length plus the
local drift error controlled by the lower bound on the reference coupling. Summing it over this interval gives
\[
 \max_i|\beta_i/\beta_n-1|
 \le \frac{H_n^{o(1)}(1+\log H_n)}{\widehat\beta_n}.
\]
In particular coupling changes are smaller than every fixed inverse
power of $H_n$, not just qualitatively $o(1)$. An interval of scales may cross any number of stages; no bound by two
stages is needed.

For the linear propagation constants the exact bound is
\begin{equation}\label{ch17:window-product}
 \log\prod_{i\text{ in window}}(2+C_{J_i})
 \le\sup_{i\text{ in window}}\varepsilon_{J_i}
                \sum_{i\text{ in window}}\log L_i
 =o(\log H_n).
\end{equation}
Thus they contribute a total factor $H_n^{o(1)}$. In particular the differentiated
curvature-response estimate in Chapter~\ref{ch11:insertions} becomes
\[
 S^{r+2}\sum_y|\Delta^r(\Psi_{n-1}\cdots\Psi_\ell)(x,y)|
 \le 2^d\prod_{i=\ell}^{n-1}R_{J_i}=H_n^{o(1)},
 \qquad S=a_\ell/a_n.
\]
The covariance decomposition telescopes for these unequal
curvature responses just as it does for equal ratios: each conditional covariance
is conjugated by the product of the subsequent curvature responses.
It uses no replacement by a common stage's quadratic form.
The finite-cutoff covariance estimate \eqref{ch11:finite-covariance}
is uniform over the blocking sequence, in particular over the unequal ratios $L_i$ of the interval of scales (Section~\ref{ch11:window-covariance}). Its stage-dependent constants,
protected-path errors, ratio-dependent jet errors and finite incidence
costs are made smaller than the requested powers of $H_n$ by choosing
sufficiently large lower bounds on the reference coupling. Sums over $O(\log H_n)$ steps lose only a factor $H_n^{o(1)}$.
At each step the linear map has its norm bound fixed before the block
ratio is chosen. The remaining nonlinear term is small once the lower bound
on the reference coupling is sufficiently large. Constants depending on the chosen ratio
therefore enter the sum of local errors. The regular and large-field source recurrences therefore retain their
strict negative powers after multiplication by the growth factor
\eqref{ch17:window-product}.

One can now repeat the proof of
Proposition~\ref{ch11:n-mark-comparison} with unequal ratios.
The initial Taylor order and support margins are chosen for the
largest $b$ and the fixed finite request. The entire incoming
polynomial in the source variables is retained. The quadratic covariance
and the polynomial $\sum_IG_Iz_I$ of free moments are common to all subsets;
the density without insertions, and hence the partition function, is unchanged.
Absorbing the larger source weight still uses
$\omega=H_n^{1/(8N_*)}$ and $I=\log H_n/(16N_*)$.
Every additional factor lost in transfer is $H_n^{o(1)}$; the margins in the
$H_n^{-3/2}$ moment error, $H_n^{-31/32}$ connected error and
$H_n^{-1}$ bounds of Lemma~\ref{ch11:restore-regular} absorb it. The chessboard bound \eqref{ch11:n-mark-root} is
therefore at most $H_n^{-c_*}$ for a fixed $c_*>0$ depending only on
the finite arity, for example $c_*=1/(32N_*)$ after increasing the lower bound on the reference coupling.
Its cell side is at most $a_nH_n^{b+5}$.
The lower bound on the partition function is imposed at the stage of the coarsest scale of the interval of scales,
with its constant at most $I$; it is not multiplied once per step.

The common localization map $T$ depends only on the insertion
and its original self covariance. It is independent of the middle
and output stopping scales. Restart on the longest interval of scales and restrict
that restarted history to the shorter ones. This gives the same $T$ for all
$b$ in the compact interval. The finite number of coefficient bounds
above is uniform in $b$: changing $b$ only changes stopping integers
and the scaling dilation $W_{H_n^b}$, whose exponents lie in a
fixed compact interval. All upper bounds used the longest interval of scales and
the largest dilation factor. No compactness over arbitrary incoming
activity arrays is invoked.

Apply this construction to
$b=\alpha_j+u-\alpha_0$, $j=0,1,2$, and
$u\in[\alpha_0,\alpha_0+1]$.
At its own normalization a single insertion has bounded chessboard seminorm.
Changing from $\alpha_j$ to $\alpha_0$ contributes a factor at most
$CH_n^{(\alpha_0-\alpha_j)D}$ by the exact scaling dilation on
the common localization map. These are precisely the size bounds used
in the two estimates above for the coincident-coordinate strips and the grid. Their exponents are
strictly negative, so the harmless factors $H_n^{o(1)}$ do not change
the conclusion. The estimates with constants independent of the number of reflected copies give the same pair bounds
in every phase. The existing argument using separated correlation pairings then gives
fixed derivatives and shell fluxes with the same diagonal subsequence and the same comparison isomorphisms $U_{i,d}$.
Thus the prescribed-radius comparison used in this chapter follows
from the construction with insertions, including intervals of scales that cross
changes of stage.
